%\section{Abstract}

Coding Theory is responsible for facilitating reliable message-passing over communication channels. In 1948, Claude Shannon proved that for every channel, there exists a family of error-correcting codes that allow reliable communication at any rate below the channel's capacity. Shannon's work is, however, non-constructive, as it provides no method to design such efficient codes. While there have been attempts over the years to reach this theoretical capacity, notably with LDPC codes (1990) and Turbo Codes (1993), these could only empirically approach the theoretical limit, and only under certain circumstances would they reach Shannon's maximum capacity. This stood for 60 years until Erdal Arikan's introduction of Polar Codes, the first clean and mathematically proven family of codes to achieve Shannon's maximum capacity over symmetric binary-input memoryless channels. \\[6pt]

This expository paper aims to briefly introduce the general Coding Theory scheme, as well as outline the key ideas in Shannon's arguments. As the first constructive example, we introduce the outstanding work of Arikan, and explore the construction behind polar codes. One concludes with ending remarks on why the discovery of polar codes presents a strong candidate for future 6G communication.